\chapter{Sets of e-variables for hypotheses generated by constraints}


\section{Sets of measures and functions}

In this section, we define and study subsets of $\mathcal{M}_s(\mathcal{X})$, the set of signed measures on $\mathcal{X}$, and $\mathcal{L}$, the set of measurable functions from $\mathcal{X}$ to $\mathbb{R}$.

\begin{definition}\label{def:integrableMeasures}
Let $\mathcal{F} \subseteq \mathcal{L}$.
We define
$
\mathcal{M}_s^{\mathcal{F}}
  = \{\mu \in \mathcal{M}_s(\mathcal{X}) \mid \forall f \in \mathcal{F}, \vert\mu\vert[\vert f \vert] < \infty\}
  \: .
$
\end{definition}


\begin{definition}\label{def:integrableFunctions}
Let $\mathcal{M} \subseteq \mathcal{M}_s(\mathcal{X})$.
We define
$
\mathcal{L}^{\mathcal{M}}
  = \{f \in \mathcal{L} \mid \forall \mu \in \mathcal{M}, \vert \mu \vert[\vert f \vert] < \infty\}
  \: .
$
\end{definition}

In \cite{larsson2025variables}, the authors consider a set of functions $\mathcal{F}$ and the set of probability measures $\mathcal{P}^{\mathcal{F}}_{\le} = \mathcal{M}_s^{\mathcal{F}} \cap \{P \in \mathcal{P}(\mathcal{X}) \mid \forall f \in \mathcal{F}, \: P[f] \le 0\}$.
Then they build $\mathcal{L}^{\mathcal{P}^{\mathcal{F}}_{\le}}$ and $\mathcal{M}_s^{\mathcal{L}^{\mathcal{P}^{\mathcal{F}}_{\le}}}$.


\begin{lemma}\label{lem:vectorSpace_integrableMeasures}
  \uses{def:integrableMeasures}
$\mathcal{M}_s^{\mathcal{F}}$ is a real vector space.
\end{lemma}

\begin{proof}

\end{proof}


\begin{lemma}\label{lem:vectorSpace_integrableFunctions}
  \uses{def:integrableFunctions}
$\mathcal{L}^{\mathcal{M}}$ is a real vector space.
\end{lemma}

\begin{proof}

\end{proof}


\begin{lemma}\label{lem:integrableMeasures_anti}
  \uses{def:integrableMeasures}
Let $\mathcal{F} \subseteq \mathcal{G} \subseteq \mathcal{L}$.
Then $\mathcal{M}_s^{\mathcal{G}} \subseteq \mathcal{M}_s^{\mathcal{F}}$.
\end{lemma}

\begin{proof}

\end{proof}


\begin{lemma}\label{lem:integrableFunctions_anti}
  \uses{def:integrableFunctions}
Let $\mathcal{M}_1 \subseteq \mathcal{M}_2 \subseteq \mathcal{M}_s(\mathcal{X})$.
Then $\mathcal{L}^{\mathcal{M}_2} \subseteq \mathcal{L}^{\mathcal{M}_1}$.
\end{lemma}

\begin{proof}

\end{proof}


\begin{lemma}\label{lem:integrableMeasures_empty}
  \uses{def:integrableMeasures}
$\mathcal{M}_s^{\emptyset} = \mathcal{M}_s(\mathcal{X})$.
\end{lemma}

\begin{proof}

\end{proof}


\begin{lemma}\label{lem:integrableFunctions_empty}
  \uses{def:integrableFunctions}
  $\mathcal{L}^{\emptyset} = \mathcal{L}$.
\end{lemma}

\begin{proof}

\end{proof}


\begin{lemma}\label{lem:bounded_subset_integrableFunctions}
  \uses{def:integrableFunctions}
  All bounded measurable functions are in $\mathcal{L}^{\mathcal{M}}$ for all $\mathcal{M} \subseteq \mathcal{M}_s(\mathcal{X})$. In particular, indicators of measurable sets are in $\mathcal{L}^{\mathcal{M}}$.
\end{lemma}

\begin{proof}
$\mathcal{M}_s(\mathcal{X})$ is the set of signed \emph{finite} measures, hence bounded functions are integrable.
\end{proof}


\begin{lemma}\label{lem:subset_integrableFunctions}
  \uses{def:integrableMeasures, def:integrableFunctions}
$\mathcal{F} \subseteq \mathcal{L}^{\mathcal{M}_s^{\mathcal{F}}}$ and $\mathcal{M} \subseteq \mathcal{M}_s^{\mathcal{L}^{\mathcal{M}}}$.
\end{lemma}

\begin{proof}
Let $f \in \mathcal{F} \subseteq \mathcal{L}$.
\begin{align*}
  f \in \mathcal{L}^{\mathcal{M}_s^{\mathcal{F}}}
  &\iff \forall \mu \in \mathcal{M}_s^{\mathcal{F}}, \: \vert \mu \vert[\vert f \vert] < \infty
  \\
  &\iff \forall \mu \in \mathcal{M}_s(\mathcal{X}), \: (\forall f' \in \mathcal{F}, \: \vert \mu \vert[\vert f' \vert] < \infty) \implies \vert \mu \vert[\vert f \vert] < \infty
\end{align*}
The last assertion is true because we can take $f' = f$.

The second inclusion is proved similarly.
\end{proof}


\begin{lemma}\label{lem:stable_integrableFunctions}
  \uses{def:integrableMeasures, def:integrableFunctions}
$\mathcal{M}_s^{\mathcal{L}^{\mathcal{M}_s^{\mathcal{F}}}} = \mathcal{M}_s^{\mathcal{F}}$ and $\mathcal{L}^{\mathcal{M}_s^{\mathcal{L}^{\mathcal{M}}}} = \mathcal{L}^{\mathcal{M}}$.
\end{lemma}

\begin{proof}
  \uses{lem:subset_integrableFunctions}
\emph{First equality}.
Since $\mathcal{F} \subseteq \mathcal{L}^{\mathcal{M}^{\mathcal{F}}}$ (Lemma~\ref{lem:subset_integrableFunctions}), we have $\mathcal{M}_s^{\mathcal{L}^{\mathcal{M}_s^{\mathcal{F}}}} \subseteq \mathcal{M}_s^{\mathcal{F}}$.
The other inclusion is Lemma~\ref{lem:subset_integrableFunctions} applied to $\mathcal{M} = \mathcal{M}_s^{\mathcal{F}}$.

\emph{Second equality}: similar.
\end{proof}


\begin{lemma}\label{lem:integral_bilin}
  \uses{def:integrableMeasures, def:integrableFunctions}
Let $\mathcal{F} \subset \mathcal{L}$ and $\mathcal{M} \subset \mathcal{M}_s(\mathcal{X})$.
Then $(\mu, f) \mapsto \mu[f]$ is a bilinear form $\mathcal{M}_s^{\mathcal{F}} \times \mathcal{L}^{\mathcal{M}_s^{\mathcal{F}}} \to \mathbb{R}$, or $\mathcal{M}_s^{\mathcal{L}^{\mathcal{M}}} \times \mathcal{L}^{\mathcal{M}} \to \mathbb{R}$.
\end{lemma}

\begin{proof}

\end{proof}


\begin{definition}\label{def:dualPair}
  \mathlibok
  \lean{LinearMap.Nondegenerate}
  A dual pair $\langle E, F \rangle$ is a pair of real vector spaces $E$ and $F$ together with a bilinear form $\langle \cdot, \cdot \rangle : E \times F \to \mathbb{R}$ such that for all $x \in E, y \in F$,
  \begin{align*}
    \forall y' \in F, &\langle x, y' \rangle = 0 \implies x = 0
    \: , \\
    \forall x' \in E, &\langle x', y \rangle = 0 \implies y = 0
    \: .
  \end{align*}
  We say that a bilinear form is \emph{non-degenerate} if the two conditions above hold.
\end{definition}


\begin{lemma}\label{lem:dualPair_integrable}
  \uses{def:dualPair, def:integrableMeasures, def:integrableFunctions}
Let $\mathcal{M} \subset \mathcal{M}_s(\mathcal{X})$.
Then $\langle \mathcal{M}_s^{\mathcal{L}^{\mathcal{M}}}, \mathcal{L}^{\mathcal{M}} \rangle$ with the integral bilinear map is a dual pair.
\end{lemma}

\begin{proof}
  \uses{lem:integral_bilin, lem:vectorSpace_integrableMeasures, lem:vectorSpace_integrableFunctions}
  We know that the two spaces are vector spaces by Lemma~\ref{lem:vectorSpace_integrableMeasures} and Lemma~\ref{lem:vectorSpace_integrableFunctions}.
  The integral is a bilinear form by Lemma~\ref{lem:integral_bilin}.
  It remains to show the separation properties.

  TODO: easy since indicators are in $\mathcal{L}^{\mathcal{M}}$ and Dirac measures are in $\mathcal{M}_s^{\mathcal{L}^{\mathcal{M}}}$.
\end{proof}


\begin{definition}\label{def:polarSet}
  \uses{def:dualPair}
  \leanok
  \lean{LinearMap.realPolar}
Let $\langle E, F \rangle$ be a dual pair. Let $S \subseteq E$. The (real) polar set of $S$ is the set of $F$ defined as
\begin{align*}
  S^\circ
  &= \{y \in F \mid \forall x \in S, \: \langle x, y \rangle \le 1\}
  \: .
\end{align*}
The bipolar set of $S$ is the set of $E$ defined as
\begin{align*}
  S^{\circ\circ}
  &= \{x \in E \mid \forall y \in S^\circ, \: \langle x, y \rangle \le 1\}
  \: .
\end{align*}
\end{definition}


\begin{lemma}\label{lem:polarSet_union}
  \uses{def:polarSet}
$(S \cup T)^\circ = S^\circ \cap T^\circ$.
\end{lemma}

\begin{proof}

\end{proof}


\begin{lemma}\label{lem:polarSet_anti}
  \uses{def:polarSet}
If $S \subseteq T$ then $T^\circ \subseteq S^\circ$ and $S^{\circ\circ} \subseteq T^{\circ\circ}$.
\end{lemma}

\begin{proof}

\end{proof}


\begin{lemma}\label{lem:polarSet_convexCone}
  \uses{def:polarSet}
If $S \subseteq E$ is a cone, then $S^\circ$ and $S^{\circ\circ}$ are also cones and can be described by
\begin{align*}
  S^\circ
  &= \{y \in F \mid \forall x \in S, \: \langle x, y \rangle \le 0\}
  \: , \\
  S^{\circ\circ}
  &= \{x \in E \mid \forall y \in S^\circ, \: \langle x, y \rangle \le 0\}
  \: .
\end{align*}
\end{lemma}

\begin{proof}

\end{proof}


\begin{definition}\label{def:weakTopology}
  \uses{def:dualPair}
  \mathlibok
  \lean{WeakBilin.instTopologicalSpace}
Let $\langle E, F \rangle$ be a dual pair.
The weak topology $\sigma(E, F)$ on $E$ is the coarsest topology such that all the maps $x \mapsto \langle x, y \rangle$ are continuous for all $y \in F$.
\end{definition}


\begin{lemma}\label{lem:isClosed_polarSet}
  \uses{def:polarSet, def:weakTopology}
Let $\langle E, F \rangle$ be a dual pair and let $S \subseteq E$.
Then $S^\circ$ is convex, contains 0 and is closed in the weak topology $\sigma(F, E)$.
\end{lemma}

\begin{proof}

\end{proof}


\begin{theorem}\label{thm:bipolar_theorem}
  \uses{def:polarSet, def:weakTopology}
Let $\langle E, F \rangle$ be a dual pair and let $S \subseteq E$.
Then $S^{\circ\circ}$ is the $\sigma(E, F)$-closure of the convex hull of $S \cup \{0\}$.
\end{theorem}

\begin{proof}

\end{proof}

Note: for two convex cones $C, K$, we have $C + K = \cc(C \cup K)$.

Returning to \cite{larsson2025variables}, we can observe that from $\mathcal{F} \subseteq \mathcal{L}$ the authors build the set of probability measures $\mathcal{P}^{\mathcal{F}}_{\le} = \mathcal{F}^\circ \cap \mathcal{P}(\mathcal{X})$, in which the dual pair is $\langle \mathcal{L}^{\mathcal{M}_s^{\mathcal{F}}}, \mathcal{M}_s^{\mathcal{F}} \rangle$.
Then they build $\mathcal{L}^{\mathcal{P}^{\mathcal{F}}_{\le}}$ and $\mathcal{M}_s^{\mathcal{L}^{\mathcal{P}^{\mathcal{F}}_{\le}}}$ and put them in duality.
We remark that the two generated spaces don't change if we replace $\mathcal{P}^{\mathcal{F}}_{\le} = \mathcal{F}^\circ \cap \mathcal{P}(\mathcal{X})$ by $\mathcal{M}^{\mathcal{F}}_{\le} = \mathcal{F}^\circ \cap \mathcal{M}(\mathcal{X})$, which is the intersection of two convex cones, hence also a convex cone.

They also consider other cones: $\mathcal{L}_p^{\mathcal{F}}$ is the cone of $\mathcal{P}^{\mathcal{F}}_{\le}$-a.e. nonnegative functions of $\mathcal{L}$.
Or: $\mathcal{L}_p^{\mathcal{P}} \subseteq \mathcal{L}^{\mathcal{P}}$ convex cone of $\mathcal{P}$-a.e. nonnegative functions of $\mathcal{L}$.

TODO: check this:
A function $f \in \mathcal{L}$ is $\mathcal{P}$-a.e. nonnegative (with $\mathcal{P}$ a set of nonnegative measures) iff for all measurable sets $A$ and all $P \in \mathcal{P}$, $P[f \mathbb{I}_A] \ge 0$.
That is, the $\mathcal{P}$-a.e. nonnegative functions are the set $\{-P \mathbb{I}_A \mid P \in \mathcal{P}, A \text{ measurable}\}^\circ$ in $\langle \mathcal{M}_s^{\mathcal{L}^{\mathcal{P}}}, \mathcal{L}^{\mathcal{P}}\rangle$.

A measure is nonnegative iff it is nonnegative on all measurable sets: $\mu[\mathbb{I}_A] \ge 0$.
The set of nonnegative measures of $\mathcal{M}_s^{\mathcal{F}}$ is $\{- \mathbb{I}_A \mid A \text{ measurable set of } \mathcal{X}\}^\circ$ (in $\langle \mathcal{L}^{\mathcal{M}_s^{\mathcal{F}}}, \mathcal{M}_s^{\mathcal{F}}\rangle$).
We get that $\mathcal{M}^{\mathcal{F}}_{\le} = \mathcal{F}^\circ \cap \{- \mathbb{I}_A \mid A \text{ measurable set of } \mathcal{X}\}^\circ$

--------------------

From $\mathcal{F} \subseteq \mathcal{L}$, we build $\mathcal{M} = \mathcal{F}^\circ \cap \{- \mathbb{I}_A \mid A \text{ measurable set of } \mathcal{X}\}^\circ$ (in the dual pair $\langle \mathcal{L}^{\mathcal{M}_s^{\mathcal{F}}}, \mathcal{M}_s^{\mathcal{F}}\rangle$).
We can then build $\mathcal{L}^{\mathcal{M}}$, which satisfies $\mathcal{L}^{\mathcal{M}_s^{\mathcal{F}}} \subseteq \mathcal{L}^{\mathcal{M}}$ since $\mathcal{M} \subseteq \mathcal{M}_s^{\mathcal{F}}$.
Then we form $\mathcal{M}_s^{\mathcal{L}^{\mathcal{M}}}$, which satisfies $\mathcal{M}_s^{\mathcal{L}^{\mathcal{M}}} \subseteq \mathcal{M}_s^{\mathcal{L}^{\mathcal{M}_s^{\mathcal{F}}}} = \mathcal{M}_s^{\mathcal{F}}$.

$\mathcal{L}_p^{\mathcal{M}}$ is the cone of $\mathcal{M}$-a.e. nonnegative functions of $\mathcal{L}^{\mathcal{M}}$.

Claim: $(\cc(\mathcal{F}) - \mathcal{L}_p^{\mathcal{M}})^\circ = \mathcal{M}$ (in $\langle \mathcal{L}^{\mathcal{M}}, \mathcal{M}_s^{\mathcal{L}^{\mathcal{M}}}\rangle$).

\begin{align*}
  (\cc(\mathcal{F}) - \mathcal{L}_p^{\mathcal{M}})^\circ
  &= (\cc(\mathcal{F}) \cup (-\mathcal{L}_p^{\mathcal{M}}))^\circ
  \\
  &= \cc(\mathcal{F})^\circ \cap (-\mathcal{L}_p^{\mathcal{M}})^\circ
  \\
  &= \mathcal{F}^\circ \cap (-\mathcal{L}_p^{\mathcal{M}})^\circ
  \\
  &\subseteq \mathcal{F}^\circ \cap \{- \mathbb{I}_A \mid A \text{ measurable set of } \mathcal{X}\}^\circ
  \\
  &= \mathcal{M}
\end{align*}
Let $\mu \in \mathcal{M}$ and $f - g \in \cc(\mathcal{F}) - \mathcal{L}_p^{\mathcal{M}}$.
Then $\mu[f - g] \le \mu[f]$ since $g \in \mathcal{L}_p^{\mathcal{M}}$. And $\mu[f] \le 0$ since $\mathcal{M} \subseteq \mathcal{F}^\circ = \cc(\mathcal{F})^\circ$.


\begin{theorem}\label{thm:eVar_fin_generated}
  \uses{def:eVar}
Let $\Phi = \{\phi_1, \ldots, \phi_d\}$ be a set of measurable functions from $\mathcal{X}$ to $\mathbb{R}$.
Let $\mathcal{P}^\Phi = \{P \in \mathcal{P}(\mathcal{X}) \mid \forall \phi \in \Phi, \ P[\vert \phi \vert] < \infty \text{ and } P[\phi] \le 0\}$.
Then $\mathcal{E}_{\mathcal{P}^\Phi}$ is equal to the set of all measurable functions $f : \mathcal{X} \to \mathbb{R}_{+, \infty}$ such that $f$ is $\mathcal{P}^{\Phi}$-a.e. lower or equal to a function of the form $1 + \sum_{i=1}^d \lambda_i \phi_i$ for some $\lambda_i \ge 0$.
\end{theorem}

\begin{proof}
See \cite[Theorem 3.1]{larsson2025variables}.
\end{proof}
